\section[Stage 003]{Stage 003: Minimal BdG--Wall Coupling and First Pole Shifts}
\label{stage:003}

\stagefield{Part / anchor}{Part I; primary anchors \resultanchor{MTDC-T2} and \resultanchor{MTDC-T4}.}
\claimstatus{The Schur complement and pole formulas are \StatusExactClosure{} inside the declared stable-mode reduction.  Passing from the full BdG PDE to stable normal coordinates is \StatusReduced{}.}

\stagefield{Purpose}{Stage~003 turns on the first conservative matter-support dynamics.  It reduces the linearized BdG sector to stable normal coordinates, couples them to the wall through the confinement variation from Stage~001, integrates them out exactly, and obtains the first microscopic origin of wall pole shifts and grouped real \(P_2\) conservative response coefficients.}

\stagefield{Inputs}{The inputs are the Stage~001 coupling \(\delta V_{\rm conf}=-(V_{\rm wall}'/\ell_c)\eta\), the Stage~002 wall coordinates, the assumption of stable BdG normal modes with positive frequencies, and harmonic selection on an isotropic reference throat.}

\stagefield{Verification}{SymPy audit: \StageFile{scripts/moving_throat_pde_stage003_bdg_sympy_audit.py}; transcript: \StageFile{scripts/output/moving_throat_pde_stage003_bdg_sympy_audit.txt}.  Mathematica audit: \StageFile{mathematica/moving_throat_pde_stage003_bdg_mathematica_audit.wl}; transcript: \StageFile{mathematica/output/moving_throat_pde_stage003_bdg_mathematica_audit.txt}.  Shared numerical stress test with Stage~004: \StageFile{scripts/numerical/stage003_004_foundational_stress.py}, \StageFile{scripts/numerical/output/stage003_004_foundational_stress.txt}, \StageFile{mathematica/numerical/stage003_004_foundational_stress.wl}, and \StageFile{mathematica/numerical/output/stage003_004_foundational_stress.txt}.}

\subsection*{Derivation}

\paragraph{1. Stable scalar-sector reduction.}
Let
\begin{equation}
Q^A=(\delta a,\delta L),
\qquad
A\in\{a,L\},
\end{equation}
and let \(X_\alpha\) be stable scalar BdG normal coordinates with frequencies \(\varpi_{0\alpha}>0\).  The reduced scalar Lagrangian is
\begin{equation}
\label{eq:app-stage003-scalar-lagrangian}
L_{\rm red}^{(0)}=
\frac12\dot Q^TM_0\dot Q-\frac12Q^TK_0Q
+\frac12\sum_\alpha\bigl(\dot X_\alpha^2-\varpi_{0\alpha}^2X_\alpha^2\bigr)
+\sum_{A,\alpha}C_{A\alpha}Q^AX_\alpha.
\end{equation}
The matrix \(C\) consists of overlap integrals generated by \eqref{eq:app-stage001-delta-v} and the BdG mode profiles.

In frequency space,
\begin{equation}
(\varpi_{0\alpha}^2-\omega^2)X_\alpha=C_{A\alpha}Q^A,
\qquad
X_\alpha=(\varpi_{0\alpha}^2-\omega^2)^{-1}C_{A\alpha}Q^A.
\end{equation}
Substitution gives the exact scalar effective kernel
\begin{equation}
\label{eq:app-stage003-scalar-kernel}
\boxed{
D_0^{\rm eff}(\omega)
=K_0-\omega^2M_0-C(\Omega_0^2-
\omega^2I)^{-1}C^T},
\end{equation}
where
\begin{equation}
\Omega_0^2=\operatorname{diag}(\varpi_{0\alpha}^2).
\end{equation}
The low-frequency expansion is
\begin{equation}
\label{eq:app-stage003-scalar-expansion}
D_0^{\rm eff}(\omega)
=K_0^{\rm eff}-\omega^2M_0^{\rm eff}-\omega^4N_0^{\rm eff}+O(\omega^6),
\end{equation}
with
\begin{equation}
\label{eq:app-stage003-keff}
\boxed{
K_0^{\rm eff}=K_0-C\Omega_0^{-2}C^T,
\qquad
M_0^{\rm eff}=M_0+C\Omega_0^{-4}C^T,
\qquad
N_0^{\rm eff}=C\Omega_0^{-6}C^T}.
\end{equation}
Stable matter support therefore lowers static stiffness, increases effective inertia, and generates the next even moment.

\paragraph{2. Grouped real \(P_2\) reduction.}
For grouped lane \(A\in\{20,21,22\}\), let \(q_A\) be the wall amplitude and \(X_{A\alpha}\) the stable BdG support coordinates.  The reduced Lagrangian is
\begin{equation}
\label{eq:app-stage003-p2-lagrangian}
L_{\rm red}^{(2)}=
\sum_A\left[
\frac12M_A\dot q_A^2-\frac12K_Aq_A^2
+\frac12\sum_\alpha(\dot X_{A\alpha}^2-\varpi_{A\alpha}^2X_{A\alpha}^2)
+\sum_\alpha g_{A\alpha}q_AX_{A\alpha}
\right].
\end{equation}
Eliminating the support coordinates gives
\begin{equation}
\label{eq:app-stage003-p2-kernel}
\boxed{
D_A^{\rm eff}(\omega)
=K_A-M_A\omega^2-
\sum_\alpha\frac{g_{A\alpha}^2}{\varpi_{A\alpha}^2-\omega^2},
\qquad
A\in\{20,21,22\}}.
\end{equation}
At low frequency,
\begin{equation}
D_A^{\rm eff}(\omega)=d_{0A}^{\rm eff}+d_{2A}^{\rm eff}\omega^2+d_{4A}^{\rm eff}\omega^4+O(\omega^6),
\end{equation}
where
\begin{align}
\label{eq:app-stage003-d0d2d4}
d_{0A}^{\rm eff}&=K_A-\sum_\alpha\frac{g_{A\alpha}^2}{\varpi_{A\alpha}^2},\\
d_{2A}^{\rm eff}&=-\left(M_A+\sum_\alpha\frac{g_{A\alpha}^2}{\varpi_{A\alpha}^4}\right),\\
d_{4A}^{\rm eff}&=-\sum_\alpha\frac{g_{A\alpha}^2}{\varpi_{A\alpha}^6}.
\end{align}
The executable audit closes finite-mode witness cases explicitly; the displayed
\(\sum_\alpha\) formulas then follow by linear superposition because each stable
BdG mode contributes additively to the Schur complement.

\paragraph{3. One-wall-mode pole shift.}
For the two-coordinate prototype
\begin{equation}
L=\frac12M\dot q^2-\frac12Kq^2+\frac12\dot X^2-\frac12\varpi^2X^2+gqX,
\end{equation}
the dispersion relation is
\begin{equation}
\label{eq:app-stage003-dispersion}
(K-M\omega^2)(\varpi^2-\omega^2)-g^2=0.
\end{equation}
Writing \(\Omega_\eta^2=K/M\), the exact poles are
\begin{equation}
\label{eq:app-stage003-poles}
\boxed{
\omega_\pm^2=\frac{\Omega_\eta^2+\varpi^2
\pm\sqrt{(\Omega_\eta^2-\varpi^2)^2+4g^2/M}}{2}}.
\end{equation}
For \(\varpi^2>\Omega_\eta^2\) and weak coupling,
\begin{equation}
\label{eq:app-stage003-pole-shift}
\boxed{
\delta\Omega_\eta^2=-\frac{g^2/M}{\varpi^2-\Omega_\eta^2}+O(g^4),
\qquad
\delta\varpi^2=+\frac{g^2/M}{\varpi^2-\Omega_\eta^2}+O(g^4)}.
\end{equation}
This is the first explicit reduced formula for how matter support shifts a wall pole.

\paragraph{4. Grouped isotropy diagnostic.}
For any grouped triple \(x_A=(x_{20},x_{21},x_{22})\), define
\begin{equation}
\label{eq:app-stage003-trace-anomaly}
\bar x=\frac{x_{20}+2x_{21}+2x_{22}}{5},
\qquad
a_x=\frac{2x_{20}-x_{21}-x_{22}}{10},
\qquad
b_x=\frac{x_{21}-x_{22}}{2}.
\end{equation}
Applying this to \(d_{2A}^{\rm eff}\) gives grouped conservative anisotropy defects.  If
\begin{equation}
K_{20}=K_{21}=K_{22},
\quad
M_{20}=M_{21}=M_{22},
\quad
g_{20,\alpha}=g_{21,\alpha}=g_{22,\alpha},
\quad
\varpi_{20,\alpha}=\varpi_{21,\alpha}=\varpi_{22,\alpha},
\end{equation}
then all grouped \(d_{nA}^{\rm eff}\) are lane-independent and
\begin{equation}
\boxed{a_x=b_x=0}.
\end{equation}
Conversely, a nonzero defect must come from anisotropy in wall data, matter support frequencies, or wall--matter overlaps.

\stagefield{Output}{Stage~003 outputs the scalar Schur complement \eqref{eq:app-stage003-scalar-kernel}, the grouped kernels \eqref{eq:app-stage003-p2-kernel}, the low-frequency coefficients \eqref{eq:app-stage003-d0d2d4}, the exact pole shift \eqref{eq:app-stage003-pole-shift}, and the first microscopic grouped-isotropy criterion.}

\stagefield{Checks}{\begin{verificationchecklist}
\item Expanding \((\Omega^2-\omega^2)^{-1}=\Omega^{-2}+\omega^2\Omega^{-4}+\omega^4\Omega^{-6}+\cdots\) gives \eqref{eq:app-stage003-keff} and \eqref{eq:app-stage003-d0d2d4}.
\item The sign of the static correction is softening: the Schur complement subtracts \(g^2/\varpi^2\).
\item The pole shift has the expected avoided-crossing sign: the lower wall-like pole moves down when coupled to a higher stable matter mode.
\item The isotropy check follows by lane equality of all microscopic inputs.
\end{verificationchecklist}}

\stagefield{Downstream use}{Stage~004 adds Maxwell/mixed self-energies to the same operator structure.  Stage~005 converts the \(D_A\) operator moments into normalized grouped response moments.  Stage~006 collects the Stage~003 BdG moments as \(B_{A0},B_{A2},B_{A4}\).}
